What is the pricing kernel? Give examples for the BSM and Binomial Tree

Consider a payoff in the Binomial tree of the form
\( S_j = e^{-rT}\sum_{i=0}^N \binom{N}{i}q^i(1 - q)^{N-i}\mathbf{1}_{i=j} = e^{-rT} \binom{N}{j}q^j(1-q)^{N-j} \)
where \( q \) are the risk-neutral probabilities and state \( j \) stands for \( S_0d^{N-j}u^j \).
Let the physical probability to end in state j be:
\( \binom{N}{i}p^i(1 - p)^{N-i} \)
We're interested in solving the optimization problem if we should consume our capital \( S_0 \)
today (call it state \( c \)) or in the future in some of the states \(c_j,\ j \in \{0, \dots, N\} \).
We normalize \( S_0 = 1 \) and solve the optimization problem

\( \max_{c, c_1, \dots, c_j} u(c_0) + \delta \sum_{j=0}^N \binom{N}{i}p^i(1 - p)^{N-i}u(c_j)\)
with the condition
\( c + \sum_{j=0}^N e^{-rT} \binom{N}{i}q^i(1 - q)^{N-i} c_j = 1 \)
The solution via Lagrange yields
\(\frac{\partial \mathcal{L}}{\partial c} = 0\)
\(u'(c) - \lambda = 0\)
\(\frac{\partial \mathcal{L}}{\partial c_{j}} = 0 \)
\({\binom{N}{j}} p^j (1-p)^{N-j} \cdot \delta \cdot u'(c_{j}) - \lambda \cdot e^{-rT} \binom{N}{j} q^j (1-q)^{N-j} = 0\)

The corresponding \( \lambda \) is the pricing kernel

\( \delta \frac{u'(u(c_j)}{u'(c)} = e^{-rT}\frac{\binom{N}{j}q^j(1 - q)^{N-j}}{\binom{N}{j}p^j(1 - p)^{N-j}} \)

It is the Radon-Nikodym derivate \( \frac{dQ}{dP} \) where \( d\mu \) is the density of the measure \( \mu \) if it exists.

In the case of the BSM, the \( dQ, dP \) are log-normal densities (with means \( r \), \( \mu \) respectively).
One can derive that
\( \delta \frac{u'(u(c_j)}{u'(c)} = \alpha S^{\frac{\mu -r}{\sigma^2}} \)
with \( \alpha \in \mathbb{R} \) with CRRA coeff. \( \frac{\mu -r}{\sigma^2} \).


Explain the basic idea for bond pricing with the pricing kernel.

Note how it holds as usual \( \mathbb{E}_Q[e^{-rT}\pi(S_T)] = \mathbb{E}_P[\pi(S_T)e^{-rT}\frac{dQ}{dP}] \).
Set \( M_t = e^{-rt}\frac{dQ}{dP} \) then pricing with the kernel for a bond would be

\( P_0 = \mathbb{E}_P[\pi(S_T)M_T(S_T)] = e^{-rT}\mathbb{E}_P[\pi(S_T)] + \text{Cov}[\pi(S_T), M_T(S_T)]\)
where the latter summand stands for the risk adjustment depending on covariation with marginal utility.

The intution is that if an asset pays off in bad states (states with high marginal
utility, i.e. an additional dollar helps a lot), this asset insures
against bad states. Most investors like these assets which
drives prices up and returns down (recall: payoff function is
fixed).

For a given probability of default \( \rho \) one can determine the price for \( 1 \) unit cash in the BSM via
\( e^{-rT}\mathbb{P}_Q[S_T > K] = e^{-rT}\Phi(-\Phi^{-1}(\rho) - \frac{(\mu - r)}{\sigma}\sqrt{T}) \).



Define the Risks of Banks and split up the aspect of credit risk

- Credit Risk
- Market Risk
- Operational Risk

Credit risk \( \text{CR} \) is the risk that a borrower defaults and
does not honor its obligation to service debt.

\( \text{CR} = \text{PD}\cdot\text{LGD}\cdot\text{EAD} \)

PD - Probability of default (associated with Default risk)

LGD - Loss Given Default (associated with Loss risk)
Loss risk determines the loss as a fraction of the exposure in the
case of default. A \( \text{LGD} < 0 \) indicates a profit (e.g. due to penalty fees)
\( \text{LGD} > 1 \) is possible for example due to ligitation costs.
The recovery rate \( \text{RR} \) is derived as \( \text{RR} = 1 - \text{LGD} \)

EAD - Exposure at the time of default (associated with Exposure risk)
The maximum amount lost at default, this isn't always obvious as for bonds,
e.g. in case of credit cards the maximum amount varies. 
Typical observation: Financially stressed counterparts have
high liquidity needs and tend to use most of the limits.

Solution: Additional clauses in the contract that allow reduced
limits or contract renegotiation when specific events occur.
Referred to as covenants or material adverse clauses.



What are PD rankings and the default rate?

PD rankings are an ordinal, relative measure of default risk,
e.g. Moody's Aaa, (Aa1, ...) to D scale or S\&P's AAA (AA+, ...) to D.

Ratings result from an (quantitative and qualitative) analysis of public and
private information. They're financed by commission fees by issuers (incentive problem).
Typical modelling techniques are

Financial models
- Structure Models (e.g., Merton's model)
- Term Structure of Default (Hazard Rates)
- Reduced-form Models (e.g., Cox process)
Empirical data based models
Expert models

Only the Financial models are relevant to the lecture.

They're usually complemented by a cardinal measure such as the default rate.
For a time period \( t \in \{0, \dots, T\} \) and rating grade \( k \), let \( n_{k,t} \) 
be the number of borrowers and \( D_{i, t} = 1 \) if the \( i \)-th borrower defaults at \( t \), then
the default rate is
\[ \bar{D}_{k,t} = \frac{\sum_{i=1}^{n_{k,t}}D_{i,t}}{n_{k,t}} \]


Give the statistical model for the probability of default (PD) from the lecture.

Let the default probability for borrower \( i \) at \( t \) be \( \pi_{i,t} \) (compare notation to default rate anki cards)
a statistical model based on generalized linear models (GLM) is:

\( \pi_{i,t} = F(\beta_0 + \beta'x_{i,t}\)
with
- \( F \) the link function
- \( \beta_0 \) the intercept (to estimate) 
- \( \beta \) factors (to estimate), \( \beta' \) for transpose,
and where \( x_{i,t} \) are exogenous variables. Typically it's components are in the lecture
- NITA = Net Income ./ Total Asset
- TLTA = (Current Liab. + Non Current Liab.) ./ Total Asset
- CASHTA = Current Assets (Cash) ./ Total Asset
- Size = log(Total Asset)
- FAR = Fixed Asset ./ Total Asset
- AGE = age of company

Examples for the link function are:

Probit function
\( F = \Phi \) then the linear predictor \( c_{i,t} = \beta_0 + \beta'x_{i,t} \) acts as a proxy
for Merton's model \( - \frac{(\mu_i - \frac{1}{2}\sigma^2_i)t + \ln(\frac{F_{i,t}}{A_{i,t}})}{\sigma_i\sqrt{t}} \).

Logistic function
\( L(s) = \frac{e^s}{1 + e^s} \)

To estimate the \( \beta \), one uses MLE, e.g. for probit

\( \ell_{i,t}(\beta, \bar{D}_{i,t}) = \phi(c_{i,t})^{\bar{D}_{i,t}}(1 - \phi(c_{i,t}))^{1 - \bar{D}_{i,t}} \)
with default rate \( \bar{D}_{i,t} \) and normal pde \( \phi \).
Then we maximize
\( \mathcal{L}(\beta, \bar{D}_i) = \prod_{t=0}^T \ell_{i,t}(\beta, \bar{D}_{i,t})\)

Example
INSERT IMAGE FROM SLIDE 29, 30 



How do ROC and confusion matrices work in the context of statistics models for PD?


A receiver operating characteristic (ROC) is a curve, created by plotting the
true positive rate (TPR) against the false positive rate (FPR) at various
threshold settings.

A large area under the ROC-curve (AUC) also known as c-statistic can be interpreted as a good model.

Confusion matrix

INSERT IMAGE SLIDE 33
INSERT IMAGE SLIDE 35



E.g. let \( \text{D} \) be the defaulted companies and \( \text{D} \) the non-defaulted,
then for given threshold \( c \), \( \text{TP} = \{i \in D | \hat{\pi}_i > c\},\ \text{FP} = \{i \in \text{ND} | \hat{\pi}_i > c\} \).

